\section{Related Work}\label{sec:related}

\subsection{Cooperative and Multi-Robot Coverage}

Coverage planning spans continuous coverage control, cellular decompositions, graph traversal, sampling-based planning, and data-driven allocation. Continuous formulations distribute mobile sensors according to spatial density functions and provide a foundation for geometric coverage control \cite{cortes2004coverage}. Sensor-based methods instead build decompositions incrementally from local observations \cite{acar2002sensor}. Surveys by Galceran and Carreras and by Tan et al. organize the large body of offline path-planning and heuristic approaches around environment representation, completeness, and path cost \cite{galceran2013survey,tan2021review}. Although these formulations differ in geometry, the coverage event is normally attached to a robot's spatial visit or sensing footprint at a planned configuration.

For multiple robots, workload division and route interaction become central. Artificial-potential-field and integrated assignment--planning approaches construct online routes while discouraging overlap or balancing the remaining region \cite{wang2024apfcpp,wang2025macplanner}. Local-search MCPP starts from a complete solution and improves it with neighborhoods that exchange or reshape route segments \cite{tang2024local}; this baseline-plus-improvement structure is important to the present planner, although ECR-Cover uses different objectives and acceptance tests. Online on-demand MCPP assigns work when execution makes it necessary, while concurrent MCPP reduces idle time by allowing planning and motion to overlap \cite{mitra2024ondemand,mitra2025concurrent}. Priority-aware coverage shifts attention from makespan alone toward when important cells are first reached \cite{lee2026priority}. Field-of-view-aware patrol planning and distribution-matching sensor deployment further show that a sensing footprint can be more informative than a point-robot path \cite{zuo2026fom,ghimire2024stein}. Persistent coverage for random-target search and inspection of complex structures illustrate related applications in which repeated visits or three-dimensional sensing geometry matter \cite{li2025persistent,jing2020multiuav}.

ECR-Cover draws on these ideas but addresses a different boundary. Its cameras are fixed in position and create coverage through pointing, so a viewpoint is a set-cover action rather than a translational cell visit. More importantly, planned feasibility and executed observation are represented by separate sets. The method does not alter the basic combinatorial difficulty of multi-agent coverage; it supplies a retained feasible baseline, a first-observation objective, and an observation-evidence interface that can be audited after execution.

\subsection{Active and Networked Camera Systems}

Pan--tilt camera research has produced systems with rapid visual feedback, target prediction, and coordinated viewpoint control. WiseCam develops systematic intelligence for pan--tilt tracking in networked settings, while high-frame-rate pan--tilt platforms couple object detection with fast actuation \cite{e2024wisecam,jiang2021pan}. Predictive scheduling and zoom control extend this direction to multi-target capture over large scenes \cite{zhang2026capture}. These systems demonstrate that camera motion, image timing, and network organization cannot be separated in a physical implementation. Their central objective, however, is usually maintaining target visibility or acquiring detailed target imagery, not certifying the first valid observation time of every spatial cell.

Multi-robot reconstruction and task-relevant exploration provide a second connection. Centralized and distributed reconstruction systems coordinate viewpoints to acquire complementary surface information \cite{hardouin2023reconstruction}; semantic Gaussian-splatting methods use scene content to guide reconstruction or exploration \cite{zeng2025multirobot,nagami2026vista}. Robust active-perception planners such as ResiVis explicitly account for uncertainty in motion and perception quality \cite{xanthidis2024resivis}. These approaches make observation value part of the planning problem. ECR-Cover is deliberately more modest in perception semantics: it does not infer target identity, reconstruction quality, or detection probability. Its observation predicate is instead operational. A cell is counted only when a usable frame acquired after stable arrival can be associated with a measured pose whose projected field of view contains that cell in the frozen camera model.

This operational definition helps separate two questions that are often conflated. Geometric feasibility asks whether a candidate pointing direction can cover a cell under the planning map. Physical confirmation asks whether the real device produced qualifying evidence for that cell during the run. A calibrated or conservative planning map may improve the former, but it cannot substitute for the latter. The distinction is particularly important with RTSP streams, for which command acknowledgment and image freshness belong to different network paths.

\subsection{Execution Uncertainty and Resilient Allocation}

Uncertain execution has been treated through stochastic scheduling, robust optimization, online replanning, and resilient assignment. Dynamic multi-robot task allocation incorporates uncertainty and temporal constraints when deciding which robot should perform a task \cite{choudhury2022dynamic}. Resilient and energy-aware allocation maintains team function under heterogeneous capabilities and robot failures \cite{notomista2022resilient}. Large-neighborhood search provides a general mechanism for starting with a feasible route and repeatedly destroying and repairing parts of it while retaining the best verified incumbent \cite{ropke2006alns}. ECR-Cover uses a small deterministic neighborhood rather than adaptive large-neighborhood search, but shares the crucial anytime principle that an interrupted search must return a validated incumbent.

The execution-time model also relates to calibrated prediction. Conformal prediction gives finite-sample marginal coverage statements under exchangeability, and Mondrian variants condition calibration on predefined groups \cite{angelopoulos2023conformal,bostrom2021mondrian}. Here, transfers are grouped by angular range so that residuals from large motions do not dominate short transfers. This construction supplies one-sided empirical margins for end-to-end time, but the formal coverage propositions do not assume that future data are automatically exchangeable. They explicitly condition on realized task times not exceeding the calibrated upper bounds. This wording keeps the role of calibration distinct from the deterministic baseline-preserving logic.

Two gaps remain when the above areas are considered together. First, a route planner can be robust with respect to its timing model and still mark a cell from a stale frame. Second, a resilient allocator can generate a plausible replacement assignment without showing that the replacement improves the coverage trajectory relative to leaving healthy queues unchanged. ECR-Cover addresses these gaps by treating the original complete route and original residual queues as executable baselines, and by making measured-pose, fresh-frame confirmation the only operation that changes the realized coverage state.

